\documentclass[conference]{IEEEtran}
\usepackage{amsmath,amssymb,graphicx,url}
\usepackage{booktabs}
\newcommand{\indicator}{\mathbf{1}}
\begin{document}
\title{IntentSat-2S: Execution Contracts for Query-Intent Satisfaction Estimation over Product Catalogs}
\author{\IEEEauthorblockN{Author information to be completed before submission}\IEEEauthorblockA{Names, affiliations, email addresses, and ORCIDs pending}}
\maketitle
\begin{abstract}
@@FULL_ABSTRACT@@
\end{abstract}
\begin{IEEEkeywords}
semantic query processing, intent satisfaction, product search, structured output, data quality, large language models
\end{IEEEkeywords}

\section{Introduction}
An accessory can match query words without satisfying a request for the device. The fraction of satisfying candidates can support pool diagnostics and monitoring if the judgments are reliable.

Large language models (LLMs) enable natural-language predicates over product descriptions, and semantic data systems expose these through declarative operators \cite{lotus,docetl,palimpzest}. Batching shares context and reduces request count, but an incomplete response can silently change an aggregate's denominator. Even complete outputs can contain incorrect judgments; these two failures require separate accounting.

We study \emph{IntentSat-2S}, which derives a query-level specification and judges candidates against it. Structural validation permits at most one retry; unresolved judgments are failures. If a specification remains invalid, an alternative plan falls back to direct batch judgment. The contract does not certify that inferred intent or judgments are correct.

We make three contributions. First, we specify a pool-level semantic aggregate, derive sharp ratio bounds conditional on an accepted-label error budget, and give deterministic candidate- and ratio-error floors. Second, we implement and compare pointwise, indexed-batch, and compact-code plans with validation, bounded retries, and fallback, including a selectively executed disagreement policy. Third, we conduct traceable studies on ESCI and an independently sourced Wayfair workload, with post-hoc LOTUS references and a biomedical transfer check. We report accuracy and pool-ratio error together with completion and inference cost, including cases in which the two-stage plan performs worse. The artifact records the sample, prompts, raw responses, model hashes, and analysis scripts. This is an initial empirical study, not a claim that the proposed plan dominates established semantic-query engines.

\section{Background and Related Work}
\textbf{Semantic data processing.} LOTUS formalizes semantic operators and optimizations relative to a reference algorithm, including accuracy guarantees under stated assumptions \cite{lotus}. Palimpzest searches for cost-quality trade-offs in AI-powered plans \cite{palimpzest}; DocETL studies logical rewrites and evaluation for document processing \cite{docetl}. SemBench evaluates multiple semantic-query engines across scenarios, operators, and modalities \cite{sembench}. Our focus is narrower: the denominator and observable failure behavior of a candidate-pool aggregate. We compare selected LOTUS filter and cascade configurations, not its complete optimizer; our structural checks do not provide LOTUS-style semantic accuracy guarantees.

\textbf{Aggregates and oracle references.} Evergreen verifies quantified claims, including proportions, with provenance; SemCEB benchmarks model-dependent cardinality estimates \cite{evergreen,semceb}. SemViews and SemReuse certify filters against recorded or fixed oracle labels \cite{semviews,semreuse}; SemViews reports nondeterminism. Our order audit highlights reference definition but does not test their guarantees. We study a relevance-labeled fixed-pool ratio, not a general claim verifier or selectivity optimizer.

\textbf{Incomplete aggregates.} Prior work reasons about query completeness over incomplete base tables \cite{query_completeness} and bounds aggregates under constraints on missing tuples \cite{contingency}. Our simpler setting fixes the candidate pool and denominator while allowing missing or erroneous generated labels; it is not a general incomplete-database framework.

\textbf{Contracts and model behavior.} Semantic integrity constraints generalize database constraints to AI-generated results, including grounding and soundness \cite{constraints}. Our implementation checks a concrete output shape and completeness, not those broader semantic properties. Prior work documents context and order sensitivity \cite{lostmiddle}, uses permutation aggregation to improve listwise ranking \cite{permcons}, and observes reordering or omission in generated topic lists \cite{outputorder}. Order sensitivity is not a new claim here: we examine its effect on an explicitly parsed candidate aggregate. Efficient serving systems, including vLLM and SGLang, optimize multi-request execution \cite{vllm,sglang}; request count alone is consequently an incomplete cost measure.

\textbf{Product relevance.} The Shopping Queries dataset provides human-assigned Exact, Substitute, Complement, and Irrelevant (E/S/C/I) relevance categories \cite{esci}. We use a documented derived encoding \cite{derived} and map E and S to positive, C and I to negative. This is a proxy for a task-specific satisfaction predicate. A substitute label does not establish that a particular shopper would purchase the item.

\section{Problem and Output Contract}
For query $q$, let $C_q=(c_1,\ldots,c_n)$ be a finite, ordered candidate pool, $n>0$. Each $c_i$ is an indexed occurrence; repeated products remain separate positions. Let $s(q,c_i)\in\{0,1\}$ denote the task's reference predicate. The target for this \emph{particular pool} is
\begin{equation}
r_q=\frac{1}{n}\sum_{i=1}^{n}s(q,c_i).
\end{equation}
A model-based plan emits $\hat{r}_q=n^{-1}\sum_i\hat{s}_i$ only when every occurrence has a valid label $\hat{s}_i$. The pool, mapping of source labels to $s$, and any text truncation must be fixed before the estimate is interpreted. A ratio over the fixed list is not, by itself, a prevalence estimate for a catalog, a user population, or another retrieval policy.

\textbf{Structural completeness.} For a chunk of $b$ candidates, array-form indexed JSON must contain exactly $b$ objects, with distinct local IDs in $\{1,\ldots,b\}$ and Y/N labels. Missing, duplicate or out-of-range IDs, duplicate JSON keys and extra fields are rejected. The parser orders labels by local ID; the orchestrator maps them back to pool positions. IDs restart in each chunk, and every position must be covered. Compact output instead contains exactly $b$ consecutive Y/N characters in input order; its checks cannot detect positional swaps. Pointwise responses contain one Y or N after whitespace trimming. An unresolved judgment leaves the pool non-estimable; specification failures have plan-specific handling.

\textbf{Sharp conditional bounds.} Suppose $m$ accepted binary labels contain $p$ positives, with $0\leq p\leq m\leq n$, and at most $k$ of them disagree with the reference, where $k\in\{0,\ldots,m\}$. With no assumptions about the $n-m$ missing labels, the reference ratio lies in
\begin{equation}
\left[\frac{\max(0,p-k)}{n},\;\frac{n-m+\min(m,p+k)}{n}\right].
\label{eq:bounds}
\end{equation}
Both endpoints are attainable over unrestricted occurrence labels: flip up to $k$ accepted Y labels to N and assign missing labels N for the lower bound, or flip accepted N labels to Y and assign missing labels Y for the upper. For duplicates sharing a reference, group $g$ has $t_g$ occurrences, $a_g$ accepted Y and $b_g$ accepted N labels. Reference 0 or 1 contributes (errors, positives) $(a_g,0)$ or $(b_g,t_g)$. A dynamic program over $G$ groups returns exact endpoints in $O(G(k+1))$ time or reports infeasibility. Interior values can be infeasible; validation uses synthetic pools. The $k=0$ case bounds completions consistent with the model labels; if no smaller error budget is justified, setting $k=m$ yields $[0,1]$. A structural parser alone does not certify a useful interval. Our operator abstains on incomplete pools; the artifact's error budget must be justified independently. Equation~\eqref{eq:bounds} is not a statistical confidence interval.

\textbf{Stability-based error floor.} For complete label vectors $a,b$ aligned to the same candidates and a fixed binary reference $s$, let $e(a,s)=n^{-1}\sum_i\indicator[a_i\ne s_i]$ and $D=\sum_i\indicator[a_i\ne b_i]$. Each disagreement forces at least one candidate error, so
\begin{equation}
\begin{aligned}
e(a,s)+e(b,s)&\ge D/n,\\
\max\{e(a,s),e(b,s)\}&\ge \lceil D/2\rceil/n.
\end{aligned}
\label{eq:stability}
\end{equation}
Both bounds are sharp over unrestricted binary references: $s=a$ attains the sum bound; distributing the $D$ disagreements evenly attains the maximum bound. Coupled duplicate labels can make the floors unattainable, though the inequalities remain valid. Agreement does not imply correctness, and candidate errors can cancel in a ratio.

For predicted positive counts $p_a,p_b$, the worse absolute ratio error is at least $\lceil|p_a-p_b|/2\rceil/n$. For $Q$ pools of equal size $n$, let $\mathrm{MAE}_a,\mathrm{MAE}_b$ be their query-macro mean absolute ratio errors. Then
\begin{equation}
\max(\mathrm{MAE}_a,\mathrm{MAE}_b)\geq
\frac{\lceil\frac{1}{2}\sum_{q=1}^{Q}|p_{a,q}-p_{b,q}|\rceil}{Qn}.
\label{eq:ratio-floor}
\end{equation}
Both floors follow from the triangle inequality and integer counts. They are sharp over unrestricted reference counts: choose a midpoint for one pool or distribute integer errors across pools. The worse output can differ by pool; neither floor identifies it or supplies a confidence interval.

\textbf{Statistical interpretation.} Conditional on the model output, $\hat{r}_q$ is a deterministic aggregate of that output. A binomial interval is not justified merely because its elements are binary: the candidate list is fixed, model decisions can be dependent, and the sampling mechanism may not represent a target population. We instead use a query-cluster bootstrap to describe variation across the sampled, constructed pools when comparing plans.

\section{Execution Plans}
Figure~\ref{fig:operator} shows the contract. The first stage maps only the query to JSON with \texttt{product\_type}, \texttt{required}, and \texttt{preferred}. The type must be a nonempty string; the other fields are string arrays. The prompt requests only query-supported constraints, but the stage can invent requirements or omit a type. The validator checks shape, not truth.

\begin{figure*}[t]
\centering\includegraphics[width=0.98\textwidth]{figures/operator.pdf}
\caption{The evaluated contract. The direct plan bypasses specification generation. The compact two-stage plan falls back to direct judgment when specification validation fails; all plans expose unresolved label failures.}
\label{fig:operator}
\end{figure*}

The second stage receives the query, products, and any specification. Indexed JSON links judgments to IDs; compact codes follow input order. Invalid batches are reissued at most once. For an invalid specification, the original two-stage plan abstains; its fallback variant uses the direct compact prompt. Batch plans partition larger pools into chunks of at most eight; every chunk must validate before a pool ratio is emitted.

For a further validation we combine two \emph{cached}, ID-aligned keyed outputs $a,b$: if $a_i=b_i$, retain that label; otherwise use a separately obtained pointwise label $h_i$. Exchanging the fixed vectors $a,b$ leaves this combination unchanged; an unresolved required label causes abstention. As a simpler comparator, sorting an unchanged chunk by stable IDs also makes its serialized prompt independent of presentation order. Neither property establishes semantic correctness, stability under changed chunks or IDs, nor deterministic results when model calls are recomputed.

For a pool of $n$ candidates, pointwise execution makes at least $n$ model calls, direct batch at least one, and two-stage batch at least two. Retries can add calls, and a fallback may replace a failed stage. These counts are not a performance proof: a batch prompt is longer, JSON generation emits more tokens, and the inference scheduler can serve many pointwise calls concurrently. We record generated tokens, offline generation time, and a local-service timing rather than inferring speed from request counts.

\section{Experimental Method}
\subsection{Data Provenance and Pool Construction}
The public \texttt{nixiesearch/amazon-esci} conversion \cite{derived} builds US-English records from the small version of Shopping Queries. Its published conversion code maps E/S/C/I to $1/0.1/0.01/0$ and emits one row per E product, together with the query's non-E products and scores. The local test file's SHA-256 matches the published file at the recorded revision. The conversion does not retain original query or product identifiers and does not include the other E products of the query in each row.

We retain the first eligible row for each exact query string, then reservoir-sample strings without replacement. A row is eligible if query and positive text are nonempty, negative texts and scores are aligned, and score codes are recognized. Each selected row contributes its single E product and at most seven randomly selected non-E products; order is shuffled with a query-specific seed. Products are truncated to 600 characters. The constructed pool therefore has an artificial composition and can omit evidence available to the original annotator. Exact-string grouping can also merge distinct original query identifiers. The derived-sample errors concern these constructed pools, not complete retrieved pages.

We use 160 query strings for an exploratory pilot. Its results motivated a compact output format and a specification-failure fallback. Before the second run we froze a new protocol and selected 500 strings from the remaining eligible test strings, excluding every pilot string. Both studies use the same source and model; the second is a disjoint query-text validation, not an independent dataset or model replication. The 500 pools contain @@NCAND@@ candidates, with source-grade counts @@GRADES@@. Ground truth maps E/S to Y and C/I to N; these are inherited relevance judgments, not newly collected purchase-satisfaction labels.

\subsection{Original Labeled Candidate Pools}
To test the effect of the derived representation, we additionally use the original Amazon Science example and product tables at a recorded revision and verify both Git LFS hashes. We exclude the exact query strings used in both previous studies and uniformly sample 250 remaining US-English test query IDs from the small version. For each sampled ID, we retain every labeled example in that version and join its product record. The resulting pools contain 5,160 labeled candidates (mean 20.64, range 8--79), with 2,293 E, 1,745 S, 242 C, and 880 I judgments. Every selected product join succeeded. These are complete \emph{labeled lists in the filtered dataset}; unjudged or live-search results may be absent. This is another sample from the same underlying dataset, not an independent data source.

\subsection{Independent Wayfair Evaluation}
After observing the ESCI results, we add the independently sourced Wayfair WANDS dataset \cite{wands}. Its annotation guidelines define Exact as a full query match, Partial as matching the entity but missing a modifier, and Irrelevant as unrelated. Unlike the ESCI E/S mapping, we use Exact as Y and both Partial and Irrelevant as N. From the source tables we collapse repeated identical labels for the same query--product pair and exclude 14 pairs with conflicting labels. We exclude query strings appearing in the three ESCI samples, then uniformly select 160 of the remaining query IDs having at least 32 nonconflicting judgments. For each selected query, we uniformly sample 32 judged products without stratifying by grade, yielding 5,120 candidates: 781 Exact, 3,282 Partial, and 1,057 Irrelevant. These are sampled labeled pools, not full retrieved pages.

We adapt the system instruction to ask for a full match to the query's stated requirements and retain the same execution contracts and output limits. Product input is a fixed, truncated concatenation of name, class, features, and description; WANDS annotators could consult product images and pages not available to the model. We repeat the four plans with both Qwen and SmolLM2 on the same 160 pools. This follow-up uses a separate data source but was designed after observing ESCI results; the new label mapping and task instruction also change, so differences cannot be attributed to the dataset alone.

\subsection{Disjoint WANDS Holdout}
@@HOLDOUT_METHOD@@

\subsection{Disagreement-Fallback Validation}
@@FALLBACK_METHOD@@

@@SELECTIVE_METHOD@@

@@ESCI_SELECTIVE_HOLDOUT_METHOD@@

\subsection{Post-hoc Local Service Test}
@@LOCAL_SERVICE_METHOD@@

\subsection{Plans and Measurements}
We compare (P) pointwise Y/N; (J) direct indexed JSON; (TJ) two-stage indexed JSON; (C) direct compact codes; and (TC) two-stage compact codes with a direct fallback. Every plan sees the same shuffled pools. J and TJ permit one full-batch retry, P one retry per invalid pointwise answer, and C and TC one compact retry. TJ abstains after repeated invalid specification output; TC falls back. Decode temperature is zero. Maximum generated tokens are four for P, 192 for indexed JSON, 160 for specification, and 32 for compact codes. On original pools we evaluate P, J, TJ, and TC, splitting batch plans into chunks of at most eight candidates. Before this run we specify all-Y, all-N, and a simple lexical baseline that marks a candidate Y only when every query word occurs in its truncated text.

Candidate precision, recall, F1, and balanced accuracy use the study-specific binary mapping---E/S positive for ESCI and Exact positive for WANDS---and are computed from pooled candidate-level confusion counts; hence larger pools contribute more candidate decisions. For the derived validation, we added all-Y and all-N constant predictors as post-hoc sanity baselines after observing the results; balanced accuracy averages sensitivity and specificity. Pool-ratio mean absolute error (MAE) averages $|\hat r_q-r_q|$ over completed pools, weighting each query equally. This is a different weighting from candidate-level F1. We report completion over all attempted pools in each study (500 derived ESCI, 250 original ESCI, or 160 WANDS); the main paired quality comparisons use pools completed by every model plan in that study. We also label a post-hoc comparison of the two JSON plans on their common pools. For paired comparisons we construct 95\% query-cluster bootstrap intervals using 4,000 resamples; the intervals are descriptive for this sampling procedure, not guarantees of population-level satisfaction. Failure-conditioned selection can still affect the common subsets.

The primary inference uses Qwen3-4B-Instruct-2507 \cite{qwen} on one NVIDIA RTX 5090 (32 GB), bfloat16, vLLM 0.10.2, PyTorch 2.8.0 with CUDA 12.8, a 4,096-token context, eager execution, and at most 16 concurrently scheduled sequences. Offline times cover chat-template rendering and batched generation calls, including retries and fallbacks; they exclude model loading and most other Python-side assembly and parsing. Prompt-token counts reflect logical request lengths, not the number of tokens actually recomputed under prefix caching. The recorded timing is not online latency or a hardware-independent comparison. After observing the primary results, we repeat the original-pool protocol with SmolLM2-1.7B-Instruct \cite{smollm} on the same 250 inputs, using the same prompts, chunk size, output limits, and GPU settings. This follow-up changes both model family and size; chat templates and tokenizers also differ. Scripts, raw outputs, model hashes, and input snapshots are retained in the artifact.

We also add a \emph{post-hoc supervised lexical baseline}: logistic regression on hashed query, product, shared-word and coverage features, trained on 50,000 original ESCI train-split examples from 18,490 distinct query strings. We exclude all case-folded, trimmed query strings in the three evaluation samples before sampling training examples, but 357 distinct product IDs appear in both its training set and the 250-pool test sample. It uses E/S versus C/I labels for training, unlike the zero-shot LLM plans, and the same truncated product text at evaluation. We fixed the features and threshold before its first run; after that run reached the optimization-iteration limit, we increased the limit to address the warning and retained both outputs. This is a useful comparator, not a model trained on an independent dataset.

@@WANDS_SUPERVISED_METHOD@@

@@WANDS_EMBEDDING_METHOD@@

For a further post-hoc reference, LOTUS 1.2.4 executes its actual \texttt{sem\_filter} on each candidate with the same Qwen model and data over a local vLLM API. Its predicate uses the same relevance proxy, but its prompt, parsing, output limit, and online execution differ from our offline pointwise plan. We retain raw answers and LOTUS decisions. The operator's configured default is True for an unparseable answer; we count such defaults separately and do not interpret this single filter as the optimized LOTUS system.

@@LOTUS_CASCADE_METHOD@@

After observing compact-format failures, we remove the four-item example and inspect strict and delimiter-tolerant parsing. We also compare the example-removed prompt with and without vLLM guided decoding, constraining each chunk to its actual number of Y/N characters. Both use each model's saved specification or fallback, candidates and order, and the same V1 engine. The artifact records an initial V0 attempt in which the constraint was not observed to take effect. As a further post-hoc audit, we reverse candidates within each fixed chunk, reassign local identifiers by position, and map the resulting labels back to the original candidates. To probe that confound, we subsequently replay the identical original prompts and reverse each chunk while retaining the identifier attached to each candidate. These runs occur in a fixed sequence and share the same pools.

After inspecting the holdout outputs, we also reuse each model's saved specification and candidate pools in an exploratory keyed-JSON audit. Each response maps every local product ID to a Y/N value, with the keys fixed in ascending ID order by guided decoding. We compare original order, identical replay and stable-ID reversal, parsing every response by ID. This changes the prompt, output grammar and token cap relative to compact output; it is not a second confirmatory holdout experiment.

\section{Results}
\subsection{Pilot and Disjoint Validation}
\textbf{Pilot.}
On the 160-query pilot, pointwise and direct JSON each completed all pools, while two-stage JSON completed 153; seven specifications lacked a valid product type after retry. Their F1 scores were 0.747, 0.585 and 0.533 on their respective completed pools, not a paired comparison. These failures motivated compact output and fallback; ESCI prompts stayed fixed across subsequent validation and original-pool studies.

\textbf{Validation.}
@@RESULTS_TEXT@@

\begin{table*}[t]
\caption{Qwen validation on 500 constructed pools. Quality is computed on the @@COMMON@@ pools completed by every LLM plan; completion, calls, token counts, and time use all 500 attempts. The constant predictors were added post hoc. Tokens are totals (thousands). Time is offline generation time on one GPU, not online latency.}
\label{tab:main}
\centering\small
\begin{tabular}{lrrrrrrrr}
\toprule
Plan & Complete & F1 & Bal. acc. & Ratio MAE & Calls & Input tok. & Output tok. & Time (s) \\
\midrule
@@TABLE_ROWS@@
\bottomrule
\end{tabular}
\end{table*}



@@INTERPRETATION@@

\textbf{Efficiency and completeness.}
@@EFFICIENCY@@

The output-format comparison should not be interpreted as a pure formatting treatment: prompts and generated text differ, so changes in the model's semantic decisions are possible. In particular, compact output enforces only the number and alphabet of decisions. Indexed output validates the presence of each identifier, but a model can still attach a wrong label to an existing identifier. The fallback improves availability only to the extent that the direct response itself validates and should not be conflated with semantic improvement.

\subsection{Original Labeled Candidate Pools and Comparators}
@@FULL_RESULTS@@

\begin{table*}[t]
\caption{Original labeled query pools. Quality is computed on the @@FULL_COMMON@@ pools completed by every Qwen plan; completion, calls, token counts, and time cover all 250 attempted pools. The constant and exact lexical baselines are pre-specified; LOTUS and the trained classifier were added post hoc. LOTUS time is excluded because it uses an online API.}
\label{tab:full}
\centering\small
\begin{tabular}{lrrrrrrrr}
\toprule
Plan & Complete & F1 & Bal. acc. & Ratio MAE & Calls & Input tok. & Output tok. & Time (s) \\
\midrule
@@FULL_TABLE_ROWS@@
\bottomrule
\end{tabular}
\end{table*}

@@FULL_INTERPRETATION@@

\textbf{Truncation sensitivity.} @@ESCI_TRUNCATION_RESULTS@@

\begin{figure*}[t]
\centering\includegraphics[width=0.98\textwidth]{figures/quality.pdf}
\caption{Quality on the original labeled pools completed by all four Qwen plans. Dots show pooled candidate F1 and mean per-pool ratio MAE, including post-hoc LOTUS filter and supervised lexical comparators; bars are 95\% query-ID-cluster bootstrap intervals. Vertical lines show pre-specified deterministic baselines. The selected pools do not measure user-level satisfaction.}
\label{fig:quality}
\end{figure*}












\textbf{LOTUS filter.}
@@LOTUS_TEXT@@

\textbf{Supervised lexical reference.}
@@SUPERVISED_TEXT@@

\textbf{Model sensitivity.}
@@SENSITIVITY_TEXT@@

\subsection{Wayfair Results and Format Audits}
@@WANDS_RESULTS@@

\begin{table*}[t]
\caption{Qwen on sampled WANDS pools using Exact as the positive label. Quality is conditional on the @@WANDS_COMMON@@ pools completed by all four model plans; completion, calls, token counts, and offline time cover all 160 attempted pools. Constant and lexical comparators are deterministic.}
\label{tab:wands}
\centering\small
\begin{tabular}{lrrrrrrrr}
\toprule
Plan & Complete & F1 & Bal. acc. & Ratio MAE & Calls & Input tok. & Output tok. & Time (s) \\
\midrule
@@WANDS_TABLE_ROWS@@
\bottomrule
\end{tabular}
\end{table*}


@@WANDS_SENSITIVITY@@

\begin{figure*}[t]
\centering\includegraphics[width=0.98\textwidth]{figures/wands_quality.pdf}
\caption{Qwen candidate F1 and query-macro ratio MAE on the WANDS pools jointly completed by all four plans. Bars are 95\% query-cluster bootstrap intervals; vertical lines show pre-specified deterministic comparators. WANDS uses an exact-only relevance proxy and sampled judged pools.}
\label{fig:wands_quality}
\end{figure*}

\textbf{Compact prompt and parser.}
@@COMPACT_ABLATION@@

@@GUIDED_COMPARISON@@

\subsection{Order, Identifiers, and Output Formats}
\textbf{Within-chunk reversal.}
@@ORDER_RESULTS@@

\begin{figure*}[t]
\centering\includegraphics[width=0.98\textwidth]{figures/order_stability.pdf}
\caption{Post-hoc order audit of the guided compact plan: candidate-label changes and mean absolute pool-ratio changes after reversing each chunk and reassigning local IDs. Dots use common completed pools; bars are 95\% query-cluster bootstrap intervals. ESCI and WANDS use different reference mappings.}
\label{fig:order_stability}
\end{figure*}

@@ORDER_CONTROL_RESULTS@@

\begin{table*}[t]
\caption{Post-hoc order controls: percentage of candidate labels that change, with 95\% query-cluster bootstrap intervals, on the common completed pools. Replay compares identical prompts; stable-ID reversal compares replay with reversal preserving each product's ID; ID reassignment compares the two reversed prompts; prior reversal compares the original with reversal and reassignment.}
\label{tab:order_controls}
\centering\small
\begin{tabular}{lrrrr}
\toprule
Model / dataset (pools) & Replay & Stable-ID reversal & ID reassignment & Prior reversal \\
\midrule
@@ORDER_CONTROL_TABLE_ROWS@@
\bottomrule
\end{tabular}
\end{table*}

\textbf{Disjoint holdout.}
@@HOLDOUT_RESULTS@@

\begin{table*}[t]
\caption{Pre-specified WANDS holdout order test on each model's common-complete pools. Change fractions are pooled over candidates; the stable-ID reversal interval is a 95\% query-cluster bootstrap. $|\Delta r|$ is the query-macro mean absolute ratio change; F1, balanced accuracy and ratio MAE evaluate the original-order judgments against the Exact-only proxy.}
\label{tab:holdout}
\centering\small
\begin{tabular}{lrrrrrr}
\toprule
Model (pools) & Replay (\%) & Stable reverse (\%, 95\% CI) & $|\Delta r|$ & F1 & Bal. acc. & Ratio MAE \\
\midrule
@@HOLDOUT_TABLE_ROWS@@
\bottomrule
\end{tabular}
\end{table*}

\textbf{Keyed-output audit.}
@@KEYED_OUTPUT_RESULTS@@

\begin{table*}[t]
\caption{Exploratory keyed-JSON audit on the exposed WANDS holdout. Entries use pools complete under all compact and keyed conditions. Change rates are pooled candidate percentages with 95\% query-cluster bootstrap intervals; compact intervals repeat Table~\ref{tab:holdout}, while keyed intervals use a separate fixed seed. F1 and ratio MAE use keyed original-order output. Compact and keyed prompts, output grammar and token caps differ.}
\label{tab:keyed}
\centering\small
\begin{tabular}{lrrrrr}
\toprule
Model (pools) & Compact reverse (\%, 95\% CI) & Keyed replay (\%) & Keyed reverse (\%, 95\% CI) & F1 & Ratio MAE \\
\midrule
@@KEYED_OUTPUT_TABLE_ROWS@@
\bottomrule
\end{tabular}
\end{table*}

\subsection{Disagreement-Fallback Validation}
@@FALLBACK_RESULTS@@

@@SELECTIVE_RESULTS@@

@@ESCI_SELECTIVE_HOLDOUT_RESULTS@@

@@LOTUS_CASCADE_RESULTS@@

\begin{table*}[t]
\caption{Disagreement-fallback and references on two disjoint WANDS samples. Quality uses common-complete pools; samples are not directly comparable. Reference training and embedding inference are excluded from generation calls and tokens. CV selects a supervised threshold using training folds. The 96-query fallbacks are trace-derived; the 52-query scheduler sends actual requests. LOTUS cascade totals include helper, calibration and routed main calls. Tokens are logical lengths across different models, not latency.}
\label{tab:fallback}
\centering\small
\begin{tabular}{lrrrrrr}
\toprule
Plan & Complete & F1 & Bal. acc. & Ratio MAE & Gen. calls & Input tokens \\
\midrule
@@FALLBACK_TABLE_ROWS@@
\bottomrule
\end{tabular}
\end{table*}

\subsection{Local Service Results}
@@LOCAL_SERVICE_RESULTS@@

\begin{table}[t]
\caption{Local HTTP workload on the existing 52 WANDS queries. Each range spans two repetitions; completion lists both counts out of 52. Latency uses only complete queries and excludes client admission, startup and warmup.}
\label{tab:local_service}
\centering\footnotesize
\begin{tabular}{lrrr}
\toprule
Plan (active queries) & p50 (s) & p95 (s) & Complete \\
\midrule
@@LOCAL_SERVICE_ROWS@@
\bottomrule
\end{tabular}
\end{table}

\subsection{Post-hoc WANDS References}
@@WANDS_SUPERVISED_RESULTS@@

@@WANDS_EMBEDDING_RESULTS@@

\subsection{Conditional Bounds on Failed Pools}
@@BOUNDS_DIAGNOSTIC@@

\subsection{Cross-Domain Transfer Check}
@@TREC_TRANSFER@@

\section{Limitations and Research Implications}
\textbf{Biased candidate pools.} Each derived ESCI row contains one E product plus non-E products, so prevalence and errors depend on that construction. The two original-pool evaluations each sample 250 query IDs and include only labeled candidates; the later sample retains some exact product-text overlap. WANDS uses 32 randomly sampled judged products per query rather than full search pages; the holdout excludes only products sampled in the earlier WANDS study. Its exact-only label mapping differs from ESCI's E/S proxy. Neither benchmark establishes the distribution of unjudged live candidates.

\textbf{Semantic and data limitations.} Truncation, missing attributes, ambiguous queries, and the E/S and Exact/Partial boundaries can alter judgments. Relevance labels do not record purchase decisions; disagreements may reflect annotation, input, or model errors. Pretraining exposure to these public benchmarks is unknown. The query-cluster bootstrap does not repair selection bias.

\textbf{Scope of the systems comparison.} Experiments use two product datasets, one biomedical dataset, two generative models and one GPU; biomedical prompts are adapted, not separately tuned. Reversals, seeded permutations and limited replays do not isolate order effects or guarantee candidate association. WANDS follow-ups reuse a source and products; the 96-query fallback uses one model. Local-service tests omit admission and production traffic. Only one LOTUS cascade is tested, not its full optimizer, Palimpzest or production search. The host had dependency conflicts; a local-copy isolated environment passes dependency and synthetic format checks, but fresh-machine installation remains unverified. We claim no general speedup or calibrated deployment-scale performance. Future work should test unjudged candidates, more domains and models, varied pool sizes and concurrency, broader systems, and independent adjudication.

\textbf{Engineering lesson.} Complete, parseable responses need not be semantically correct. Compact representations and fallbacks expose cost and error trade-offs that must be measured.

\section{Conclusion}
We defined a query-intent ratio over a declared pool and implemented observable failures for several LLM plans. Completeness, semantic quality, and execution cost must be assessed separately: fewer requests do not by themselves establish a better plan. Further data and evaluation are required before deployment.

\section*{Acknowledgments and AI-Generated Content Disclosure}
OpenAI Codex (GPT-6) substantially drafted and revised Sections I--VIII, wrote evaluation and analysis scripts and the code for Figures 1--4, and assisted literature review and interpretation. Qwen3-4B-Instruct-2507 and SmolLM2-1.7B-Instruct generated Section VI judgments; Qwen3-4B also generated Section VII's synthetic format checks. Qwen3-Embedding-0.6B generated vectors under Section V's protocol. The authors are responsible for the correctness and originality of the text, code, figures, sources and results.

\IEEEtriggeratref{16}
\begin{thebibliography}{99}
\bibitem{lotus} L. Patel, S. Jha, M. Pan, H. Gupta, P. Asawa, C. Guestrin, and M. Zaharia, ``Semantic operators and their optimization: Enabling LLM-based data processing with accuracy guarantees in LOTUS,'' \emph{Proc. VLDB Endow.}, vol. 18, no. 11, pp. 4171--4184, 2025, doi: 10.14778/3749646.3749685.
\bibitem{docetl} S. Shankar, T. Chambers, T. Shah, A. G. Parameswaran, and E. Wu, ``DocETL: Agentic query rewriting and evaluation for complex document processing,'' \emph{Proc. VLDB Endow.}, vol. 18, no. 9, pp. 3035--3048, 2025, doi: 10.14778/3746405.3746426.
\bibitem{palimpzest} C. Liu \emph{et al.}, ``Palimpzest: Optimizing AI-powered analytics with declarative query processing,'' in \emph{Proc. CIDR}, 2025.
\bibitem{sembench} J. Lao \emph{et al.}, ``SemBench: A benchmark for semantic query processing engines,'' \emph{Proc. VLDB Endow.}, vol. 19, no. 8, pp. 1754--1767, 2026, doi: 10.14778/3811243.3811249.
\bibitem{evergreen} A. W. Lee, B. Han, S. Sen, S. Yeom, U. Cetintemel, and A. Datta, ``Evergreen: Efficient claim verification for semantic aggregates,'' arXiv:2604.26180, 2026.
\bibitem{semceb} A. Zimmerer, C. K\"uhn, Y. Li, M. Stoian, R. Borovica-Gajic, and A. Kipf, ``SemCEB: A cardinality estimation benchmark for semantic operators,'' arXiv:2606.23081, 2026.
\bibitem{semviews} O. Hassanzadeh and D. Subramanian, ``Semantic views: Certified reuse of LLM labels across semantic queries,'' arXiv:2610.12768, 2026.
\bibitem{semreuse} P. T. Tran-Truong, T. Kieu, X.-B. Le, S. X. Ha, D. Kieu, B. Yang, and C. S. Jensen, ``SemReuse: Certified reuse of cached natural-language views for semantic operators,'' arXiv:2610.16035, 2026.
\bibitem{query_completeness} S. Razniewski and W. Nutt, ``Completeness of queries over incomplete databases,'' \emph{Proc. VLDB Endow.}, vol. 4, no. 11, pp. 749--760, 2011, doi: 10.14778/3402707.3402715.
\bibitem{contingency} X. Liang, Z. Shang, S. Krishnan, A. J. Elmore, and M. J. Franklin, ``Fast and reliable missing data contingency analysis with predicate-constraints,'' in \emph{Proc. SIGMOD}, 2020, pp. 285--295, doi: 10.1145/3318464.3389785.
\bibitem{constraints} A. W. Lee \emph{et al.}, ``Semantic integrity constraints: Declarative guardrails for AI-augmented data processing systems,'' \emph{Proc. VLDB Endow.}, vol. 18, no. 11, pp. 4073--4080, 2025, doi: 10.14778/3749646.3749677.
\bibitem{lostmiddle} N. F. Liu \emph{et al.}, ``Lost in the middle: How language models use long contexts,'' \emph{Trans. Assoc. Comput. Linguistics}, vol. 12, pp. 157--173, 2024, doi: 10.1162/tacl\_a\_00638.
\bibitem{permcons} R. Tang, C. Zhang, X. Ma, J. Lin, and F. Ture, ``Found in the middle: Permutation self-consistency improves listwise ranking in large language models,'' in \emph{Proc. NAACL}, 2024, pp. 2327--2340, doi: 10.18653/v1/2024.naacl-long.129.
\bibitem{outputorder} P. Schilcher \emph{et al.}, ``Characterizing positional bias in large language models: A multi-model evaluation of prompt order effects,'' in \emph{Findings of EMNLP}, 2025, pp. 20643--20664, doi: 10.18653/v1/2025.findings-emnlp.1124.
\bibitem{vllm} W. Kwon \emph{et al.}, ``Efficient memory management for large language model serving with PagedAttention,'' in \emph{Proc. SOSP}, 2023, pp. 611--626, doi: 10.1145/3600006.3613165.
\bibitem{sglang} L. Zheng \emph{et al.}, ``SGLang: Efficient execution of structured language model programs,'' \emph{Adv. Neural Inf. Process. Syst.}, vol. 37, pp. 62557--62583, 2024, doi: 10.52202/079017-2000.
\bibitem{esci} C. K. Reddy \emph{et al.}, ``Shopping Queries Dataset: A large-scale ESCI benchmark for improving product search,'' arXiv:2206.06588, 2022.
\bibitem{derived} Nixiesearch, ``Amazon ESCI dataset in nixietune format,'' revision c092c898, 2024. [Online]. Available: \url{https://huggingface.co/datasets/nixiesearch/amazon-esci}
\bibitem{wands} Y. Chen, S. Liu, Z. Liu, W. Sun, L. Baltrunas, and B. Schroeder, ``WANDS: Dataset for product search relevance assessment,'' in \emph{Proc. ECIR}, 2022, pp. 128--141, doi: 10.1007/978-3-030-99736-6\_9.
\bibitem{pigeonhole} A. B. Owen, ``The pigeonhole bootstrap,'' \emph{Ann. Appl. Stat.}, vol. 1, no. 2, pp. 386--411, 2007, doi: 10.1214/07-AOAS122.
\bibitem{qwen} Qwen, ``Qwen3-4B-Instruct-2507,'' model card, 2025. [Online]. Available: \url{https://huggingface.co/Qwen/Qwen3-4B-Instruct-2507}
\bibitem{smollm} Hugging Face, ``SmolLM2-1.7B-Instruct,'' model card, 2024. [Online]. Available: \url{https://huggingface.co/HuggingFaceTB/SmolLM2-1.7B-Instruct}
\bibitem{qwen_embed} Qwen, ``Qwen3-Embedding-0.6B,'' model card, 2025. [Online]. Available: \url{https://huggingface.co/Qwen/Qwen3-Embedding-0.6B}
\bibitem{treccovid} K. Roberts \emph{et al.}, ``Searching for scientific evidence in a pandemic: An overview of TREC-COVID,'' \emph{J. Biomed. Inform.}, vol. 121, p. 103865, 2021, doi: 10.1016/j.jbi.2021.103865.
\end{thebibliography}
\end{document}
